%%%%%%%%%%%%%%%%%%%%%%%%%%%%%%%%%%%%%%%%%%%%%%%%%%%%%%%%%%%%%%%%%%%%%
%
%   Formal Derivation Whitepaper Template
%   A clean, standard format for mathematical and scientific derivations.
%
%%%%%%%%%%%%%%%%%%%%%%%%%%%%%%%%%%%%%%%%%%%%%%%%%%%%%%%%%%%%%%%%%%%%%

\documentclass[11pt, letterpaper]{article}

% PACKAGES
\usepackage[margin=1in]{geometry} % Set page margins
\usepackage{amsmath, amssymb}     % For advanced math typesetting
\usepackage{graphicx}             % For including figures
\usepackage{hyperref}             % For clickable links and references
\usepackage{lipsum}               % For generating placeholder text

% --- DOCUMENT INFORMATION ---
\title{Formal Derivation of [Your Topic Here]}
\author{Justin K. Lietz \\[0.2em] \textit{Founder of Neuroca, Inc.}}
\date{\today} % Or \date{August 24, 2025} for a specific date

% --- BEGIN DOCUMENT ---
\begin{document}

\maketitle

% --- ABSTRACT ---
\begin{abstract}
    This document presents a formal, step-by-step derivation of [briefly state the result]. Starting from the foundational principles of [mention starting principles], we systematically develop the continuum equations of motion. The analysis clarifies the origin of key terms and ensures dimensional consistency throughout the derivation process. The final result provides a rigorous mathematical foundation for [mention the application or significance].
    \vspace{1cm} % Adds some space after the abstract
\end{abstract}


% --- SECTIONS ---

\section{Introduction}
\lipsum[1] % Placeholder text. Replace with your introduction.
The primary objective is to bridge the gap between the discrete formulation and its macroscopic continuum limit. We will begin by outlining the initial assumptions and definitions.

\section{Preliminaries and Definitions}
Before proceeding with the main derivation, we must establish the foundational concepts and notation.

\subsection{The Discrete Dynamical Law}
The system originates from a local update rule governing the state $W_i(t)$ of a discrete node $i$. The fundamental law is given by:
\begin{equation}
    \frac{\Delta W_i}{\Delta t} = (\alpha - \beta)W_i - \alpha W_i^2 + J \sum_{j\in \mathrm{nbr}(i)} (W_j - W_i)
    \label{eq:discrete_law}
\end{equation}
Here, $\alpha$ and $\beta$ represent on-site parameters, and $J$ is the coupling constant.

\subsection{Continuum Field Postulate}
We postulate a continuous scalar field $\phi(\mathbf{x}, t)$ that represents the coarse-grained average of the discrete states $W_i(t)$. The mapping is formally expressed as:
$$
    W_i(t) \quad \to \quad \phi(\mathbf{x}_i, t)
$$
where $\mathbf{x}_i$ is the spatial position of node $i$.

\section{Main Derivation}
This section contains the core mathematical steps of the derivation.

\subsection{Temporal Approximation}
In the limit of an infinitesimally small time step, $\Delta t \to 0$, the left-hand side of Equation \ref{eq:discrete_law} becomes a partial time derivative:
\begin{equation}
    \lim_{\Delta t \to 0} \frac{\Delta W_i}{\Delta t} = \frac{\partial \phi}{\partial t}
\end{equation}

\subsection{Spatial Approximation (Discrete Laplacian)}
The interaction term is a discrete representation of the Laplacian operator. For a regular lattice with spacing $a$, a Taylor series expansion reveals this relationship. The derivation proceeds as follows:
\begin{align}
    \sum_{j\in \mathrm{nbr}(i)} (W_j - W_i) &\approx \sum_{j\in \mathrm{nbr}(i)} \left( \phi(\mathbf{x}_i + \mathbf{\delta}_j) - \phi(\mathbf{x}_i) \right) \\
    &\approx C_L a^2 \nabla^2 \phi
\end{align}
where $C_L$ is a lattice-dependent constant and we define the diffusion coefficient $D = J C_L a^2$.

\subsection{Assembling the Final Equation}
By substituting the temporal and spatial approximations into the discrete law, we arrive at the final continuum partial differential equation
\begin{equation}
    \partial_t \phi = D\, \nabla^{2}\phi + r\, \phi - u\, \phi^{2}
    \label{eq:final_pde}
\end{equation}
where we have defined $r = \alpha - \beta$ and $u = \alpha$ for clarity.

\section{Analysis and Discussion}
\lipsum[2] % Placeholder text. Replace with your analysis of the result.
Equation \ref{eq:final_pde} is a well-known Reaction-Diffusion equation. The stability of its fixed points can be analyzed by setting the time and spatial derivatives to zero.

\section{Conclusion}
\lipsum[3] % Placeholder text. Replace with your conclusion.
This document has formally derived the continuum equation of motion from the model's discrete foundations, providing a rigorous link between the microscopic rules and the macroscopic behavior.

\vspace{2cm} % Add vertical space before the signature

% --- SIGNATURE BLOCK ---
\noindent % Prevents indentation
\begin{minipage}{0.5\linewidth}
    \hrulefill \\ % The signature line
    \vspace{0.5em}
    \noindent \textbf{Justin K. Lietz} \\
    \textit{Founder of Neuroca, Inc.}
\end{minipage}


\end{document}